\begin{textsample}{NEG Dim 5 – ba – Score -3.00 – ba/ba\_inf\_1.txt}  \label{ex:f5_neg_008}
7.1.
Introdução No fluxo e na dinâmica da implementação da Base Nacional Comum Curricular para a Educação Infantil, efetiva -se e explicita -se aqui o Referencial Curricular para a Educação Infantil.
Faz -se necessário afirmar, como princípio fundante, que a BNCC não pode se configurar enquanto norma que visa a uma mera coesão curricular no currículo do Estado, do Município ou da Escola.
Trata -se de uma normatização que orienta, referencia, mas que não deve se configurar como uma prescrição totalizante, até porque, nos processos de democratização das realizações curriculares, não há nem deverá haver autoridades curriculares únicas se levarmos em conta os contextos socioeducacionais de apropriação curricular.
Na busca por uma “ coesão ” pretendida pela BNCC, como política de currículo, com autonomia e responsabilidade socioeducacional, o DCRB movimenta -se numa apropriação crítico-contextualizada de suas orientações.
Faz -se necessário afirmar que este documento não pode ser tomado como uma prescrição curricular, mas um Referencial que, para ser pertinente e relevante, deve ser estudado e debatido, para depois desse processo ser apropriado pelos educadores da Educação Infantil nos contextos municipais, nas escolas e salas de aula.
Levando em conta a extensão e a diversidade territorial, socioeconômica e cultural do Estado da Bahia, assim como suas complexas e intensas necessidades educacionais, com implicações diretas de como emerge a condição concreta das infâncias baianas e o que exigem historicamente em termos das necessidades socioeducacionais para uma formação qualificada, esse Referencial se configura dentro de uma perspectiva multirreferencial, na medida em que se inspira na compreensão de que a criança não é uma invariante existencial, social e/ou cultural.
Dessa forma é que, para fins curriculares, pensamos no bem comum socialmente referenciado e na emergência e acolhimento da diferença como realidades com as quais devemos trabalhar dispositivos curriculares e processos formativos e aprendermos mutuamente no convívio socioeducacional, em face da sua riqueza e importância formativo-experiencial.
Não abrindo mão dos processos de autonomização e de responsabilização da escola no que concerne à relação com as normas e à ineliminável e complexa dinâmica curricular e seus “ atos de currículo ” ( MACEDO, 2013 ; 2016 ), fundamental é pensarmos nas crianças em formação como atores/atrizes e autores/autoras curriculantes.
Com suas experiências, ações e protagonismos aprendentes, elas podem acrescentar cotidianamente ao currículo, assim como os grupos e instituições com os quais interagem, tomando -os como analisadores e propositores de processos decisórios da gestão curricular.
É importante, neste processo, a intercriticidade das suas contribuições à experiência e à dinâmica curricular.
Nesses termos, são instituintes de atos de currículo.
É assim que autoformação ( formar -se ), heteroformação ( formar -se com o outro ) e metaformação ( reflexão voltada para a análise \textbf{valorada} do próprio processo formativo ) passam a se configurar como realizações importantes para a qualificação da formação da criança.
Nesses termos, pedagogias ativas, ou seja, construcionistas, social e culturalmente referenciadas, são apropriadas em favor de uma didática coerente com essas perspectivas curriculares e as aprendizagens aí requeridas.
É assim que cuidado e aprendizagem qualificada e a aprendizagem qualificada pelo cuidado, bem como interações singularizadas, tomando a criança como ser de desejo, intenção e pensamento, vivendo de forma entretecida de as experiências corporais, cognitivas e socioculturais, assim como inserida e atuante nas complexas dinâmicas das sociedades em que vivem sejam consideradas sujeitos de fato e de direitos.
É fundamental imaginar um currículo para Educação Infantil inserido reflexiva e criticamente nas intensas, enlarguecidas e dinâmicas culturais da contemporaneidade.
Currículos contemporâneos se configuram pela vivência e apropriação de saberes e atividades que se complementam, que vivem reflexiva e intensamente a experiência e o acontecimento e que, no presente, se nos apresentam com potencialidades para qualificação da formação.
Num mundo cibercultural, por exemplo, entre outros mundos culturais, onde as crianças os vivenciam com intenso protagonismo, numa dinâmica, muitas vezes, ambivalente no que se refere às questões éticas, afetivas e cognitivas.
Na mesma medida dessa intensidade experiencial, não há como um currículo e seus atos deixarem de criar as condições para se configurar processos de aprendizagem qualificados, implicando experiências e experimentações aprendentes, nas quais as aprendizagens deverão articular saberes curriculares vinculados aos campos dos saberes sistematizados, da experiência, da ética, da estética, da política, da cultura e da espiritualidade.
Com isso, pautas vinculadas à justiça social, como questões de desigualdade socioeconômica, de gênero, étnico-raciais, geracionais e de orientação sexual são transversais à formação da criança.
Trata -se aqui de uma formação cidadã qualificada pela interseccionalidade da formação.
Ademais e de forma realçada, qualquer experiência curricular e seus atos de currículo devem ter como preocupação contínua reflexões aprofundadas sobre qual a concepção de infância que norteia o currículo, por se tratar de uma construção filosófica, sociocultural e política fundamental para se qualificar um currículo de Educação Infantil.

% matched lemmas: valorar
\end{textsample}
